% Options for packages loaded elsewhere
\PassOptionsToPackage{unicode}{hyperref}
\PassOptionsToPackage{hyphens}{url}
%
\documentclass[
  10pt,
]{article}
\usepackage{amsmath,amssymb}
\usepackage{iftex}
\ifPDFTeX
  \usepackage[T1]{fontenc}
  \usepackage[utf8]{inputenc}
  \usepackage{textcomp} % provide euro and other symbols
\else % if luatex or xetex
  \usepackage{unicode-math} % this also loads fontspec
  \defaultfontfeatures{Scale=MatchLowercase}
  \defaultfontfeatures[\rmfamily]{Ligatures=TeX,Scale=1}
\fi
\usepackage{lmodern}
\ifPDFTeX\else
  % xetex/luatex font selection
\fi
% Use upquote if available, for straight quotes in verbatim environments
\IfFileExists{upquote.sty}{\usepackage{upquote}}{}
\IfFileExists{microtype.sty}{% use microtype if available
  \usepackage[]{microtype}
  \UseMicrotypeSet[protrusion]{basicmath} % disable protrusion for tt fonts
}{}
\makeatletter
\@ifundefined{KOMAClassName}{% if non-KOMA class
  \IfFileExists{parskip.sty}{%
    \usepackage{parskip}
  }{% else
    \setlength{\parindent}{0pt}
    \setlength{\parskip}{6pt plus 2pt minus 1pt}}
}{% if KOMA class
  \KOMAoptions{parskip=half}}
\makeatother
\usepackage{xcolor}
\usepackage[margin=0.85in]{geometry}
\usepackage{longtable,booktabs,array}
\usepackage{calc} % for calculating minipage widths
% Correct order of tables after \paragraph or \subparagraph
\usepackage{etoolbox}
\makeatletter
\patchcmd\longtable{\par}{\if@noskipsec\mbox{}\fi\par}{}{}
\makeatother
% Allow footnotes in longtable head/foot
\IfFileExists{footnotehyper.sty}{\usepackage{footnotehyper}}{\usepackage{footnote}}
\makesavenoteenv{longtable}
\setlength{\emergencystretch}{3em} % prevent overfull lines
\providecommand{\tightlist}{%
  \setlength{\itemsep}{0pt}\setlength{\parskip}{0pt}}
\setcounter{secnumdepth}{-\maxdimen} % remove section numbering
\usepackage{amsmath,amssymb,mathtools,booktabs,longtable}
\usepackage{microtype}
\usepackage[hidelinks]{hyperref}
\ifLuaTeX
  \usepackage{selnolig}  % disable illegal ligatures
\fi
\usepackage{bookmark}
\IfFileExists{xurl.sty}{\usepackage{xurl}}{} % add URL line breaks if available
\urlstyle{same}
\hypersetup{
  pdftitle={WHestBench Complete Proof Package},
  pdfauthor={Research synthesis, 2026-07-30},
  hidelinks,
  pdfcreator={LaTeX via pandoc}}

\title{WHestBench Complete Proof Package}
\usepackage{etoolbox}
\makeatletter
\providecommand{\subtitle}[1]{% add subtitle to \maketitle
  \apptocmd{\@title}{\par {\large #1 \par}}{}{}
}
\makeatother
\subtitle{Geometric optimality, correction observability, exact control
means, and frozen empirical certificates}
\author{Research synthesis, 2026-07-30}
\date{2026-07-30}

\begin{document}
\maketitle

\section{1. Executive theorem map}\label{executive-theorem-map}

This document gives the strongest rigorous synthesis supported by the
WHestBench proof and experiment artifacts. It deliberately distinguishes
four statuses:

\begin{enumerate}
\def\labelenumi{\arabic{enumi}.}
\tightlist
\item
  \textbf{Analytically proved:} the proof is included below.
\item
  \textbf{Computer-assisted certified dependency:} a separate
  interval/exact-arithmetic package supplies the proof; this document
  states the theorem, its trust base, and the logical consequences.
\item
  \textbf{Finite frozen certificate:} a deterministic verifier
  recomputes a claim from the saved binary artifacts.
\item
  \textbf{Empirical only:} the result motivates a theorem or closes a
  tested implementation, but does not imply population-level
  impossibility.
\end{enumerate}

The central conclusion is not that no competition winner exists. The
valid conclusion is a \textbf{two-boundary theorem}:

\begin{itemize}
\tightlist
\item
  \textbf{Geometric boundary.} In the specified infinite-width, static,
  network-independent cubature classes, complete Kerdock is certified
  near-optimal, and is exactly optimal inside the fixed Kerdock-line
  universe even when arbitrary real line weights are allowed.
\item
  \textbf{Information boundary.} For any explicitly specified runtime
  information class, the best possible adaptive correction is the
  Hilbert projection of the baseline error onto that information class.
  Same-design centered information, Haar-orientation-invariant
  information, or any feature class with a phase-flip symmetry can have
  exactly zero correction value even when a large per-network oracle
  exists.
\end{itemize}

A material improvement must cross at least one boundary by introducing
new geometry, new absolute-phase information, a transformed residual
kernel, finite-width-specific structure, or a computationally richer
white-box procedure.

\subsection{1.1 Claim-status table}\label{claim-status-table}

\begin{longtable}[]{@{}
  >{\raggedright\arraybackslash}p{(\columnwidth - 6\tabcolsep) * \real{0.2500}}
  >{\raggedright\arraybackslash}p{(\columnwidth - 6\tabcolsep) * \real{0.2500}}
  >{\raggedright\arraybackslash}p{(\columnwidth - 6\tabcolsep) * \real{0.2500}}
  >{\raggedright\arraybackslash}p{(\columnwidth - 6\tabcolsep) * \real{0.2500}}@{}}
\toprule\noalign{}
\begin{minipage}[b]{\linewidth}\raggedright
ID
\end{minipage} & \begin{minipage}[b]{\linewidth}\raggedright
Claim
\end{minipage} & \begin{minipage}[b]{\linewidth}\raggedright
Status
\end{minipage} & \begin{minipage}[b]{\linewidth}\raggedright
Scope
\end{minipage} \\
\midrule\noalign{}
\endhead
\bottomrule\noalign{}
\endlastfoot
G1 & Kerdock is within \texttt{0.02336550102948\%} of the static
nonnegative optimum & Computer-assisted certified & Infinite-width
depth-32 ReLU kernel, \texttt{d=256}, at most 66,048 nodes, nonnegative
mass-one weights, network-independent rule \\
G2 & Tightened auxiliary-LP comparison gives
\texttt{0.023324172950039\%} relative excess & Computer-assisted
certified & Same limiting kernel; all-degree admissible auxiliary LP \\
G3 & Complete bases plus at most one partial basis exactly optimize
arbitrary real weights on the fixed Kerdock-line universe & Analytically
proved & 33,024 symmetrized lines; static linear rules \\
G4 & Any improvement exceeding \texttt{0.02331873404818\%} must leave
the tightened static class & Analytically proved corollary & Same as
G2 \\
I1 & Runtime information value equals squared conditional-projection
norm & Analytically proved & Any square-integrable Hilbert-valued
error \\
I2 & Optimal adaptive coefficients in a finite correction dictionary
satisfy conditional normal equations & Analytically proved &
Finite-dimensional dictionary, square-integrable features \\
I3 & Phase-flip symmetry forces zero correction value & Analytically
proved & Explicit measure-preserving involution and invariant
information \\
I4 & Haar-random relative orientation makes every orientation-blind
correction useless & Analytically proved & Independent Haar orientation
and mass-one cubature \\
I5 & Same-design common bias is minimax non-identifiable & Analytically
proved & Equal-loading observation model \\
I6 & Full weights defeat a pure information-theoretic impossibility
theorem & Analytically proved & Target/error deterministic from complete
weights \\
R1 & Exact correction-risk, shrinkage, and downstream replacement
formulas & Analytically proved & Hilbert-space scoring model \\
R2 & ReLU nonlinear remainder is supported on gate crossings and has a
cubic second-moment bound & Analytically proved & Explicit density
assumption near zero \\
C1 & Controls in the quadrature exactness space are pathwise no-ops,
even with adaptive global parameters & Analytically proved & Uniformly
annihilated control family \\
C2 & Symmetrized spherical Poisson controls have exact mean one and
nonzero all-even harmonics & Analytically proved & Sphere,
\texttt{\textbar{}r\textbar{}\textless{}1} \\
C3 & Biased projected-ReLU controls have a closed-form exact spherical
mean & Analytically proved & Uniform fixed-radius sphere \\
E1 & Ordinary independent replication is score-neutral under linear
MSE-times-compute accounting & Analytically proved & Equal costs,
independent errors; correlated extension included \\
F1 & Kernel near-optimality transfers only under a quantitative uniform
finite-width kernel bound & Analytically proved & Bounded total
variation of rules \\
F2 & Analytic-plus-residual methods require recertification using the
residual kernel & Analytically proved & Exact integral of the analytic
component \\
X1 & Frozen Poisson, projected-ReLU, and signed-probe candidates fail
their stated gates & Finite frozen certificate & Saved IEEE-754 arrays
and frozen scripts only \\
X2 & A universal no-winning-estimator theorem follows from the
experiments & \textbf{False / not proved} & Countered by the
full-information theorem and open estimator classes \\
\end{longtable}

\section{2. Mathematical setup}\label{mathematical-setup}

Let

\begin{itemize}
\tightlist
\item
  \((\Omega,\mathcal F,\mathbb P)\) be the randomness space;
\item
  \(H\) be the real Hilbert space of scored output coordinates;
\item
  \(I(f)\) be the true spherical integral of an integrand \(f\);
\item
  \(Q(f)=\sum_{i=1}^m w_i f(x_i)\) be a cubature rule with
  \(\sum_i w_i=1\);
\item
  \(\widehat y_0\) be the protected baseline estimate;
\item
  \(y\) be the true target;
\item
  \(e=\widehat y_0-y\in L^2(\Omega;H)\) be baseline error.
\end{itemize}

A correction \(c\) is applied as \(\widehat y=\widehat y_0-c\), so its
risk is

\[
R(c)=\mathbb E\|e-c\|_H^2.
\]

The competition-relevant empirical ratios use

\[
\text{raw ratio}=\frac{\text{candidate MSE}}{\text{baseline MSE}},
\qquad
\text{adjusted ratio}=\text{raw ratio}\times\text{compute ratio}.
\]

The cited \texttt{4.34x} improvement target corresponds to an adjusted
ratio at most

\[
\frac1{4.34}=0.2304147465\ldots.
\]

\section{3. Geometric boundary}\label{geometric-boundary}

\subsection{3.1 Imported certified result: static nonnegative
near-optimality}\label{imported-certified-result-static-nonnegative-near-optimality}

For the normalized infinite-width depth-32 ReLU kernel in dimension
\(256\), the T22 release proves a one-sided bound of the form

\[
R_K\le (1+\delta_{22})R_*,
\qquad
\delta_{22}=0.0002336550102948\ldots,
\]

where \(R_K\) is complete-Kerdock ensemble MSE and \(R_*\) is the
infimum over static, network-independent, nonnegative mass-one cubature
rules using at most 66,048 spherical nodes.

The release is a computer-assisted proof, not a proof-assistant
formalization. Its documented trust base includes exact-rational
witnesses, directed rounding, 1,421 certified pointwise intervals, a
directed kernel-mean enclosure, one-sided ratio logic, and a 59-entry
manifest. The stale two-sided artifact is excluded.

The later all-degree auxiliary-LP certificate gives the stronger
comparison

\[
R_K\le (1+\delta_{30})R_*,
\qquad
\delta_{30}=0.00023324172950039,
\]

or \texttt{0.023324172950039\%} relative excess. Its separate proof uses
the degree-five Hermite interpolant at the roots of

\[
22102t^3+21930t^2-87t-85=0,
\]

strict positivity of \(K_{32}^{(6)}\), strict negativity of every unused
reduced cost above degree five, positive nonconstant Gegenbauer
coefficients, and exact primal-dual equality.

These are imported certified dependencies. The bundled verifier in this
package does not rerun their interval engines.

\subsection{3.2 Exact structural escape
threshold}\label{exact-structural-escape-threshold}

\textbf{Corollary 3.1.} Every estimator in the tightened static
nonnegative class satisfies

\[
R\ge \frac{R_K}{1+\delta_{30}}.
\]

Hence its fractional improvement over Kerdock is at most

\[
1-\frac1{1+\delta_{30}}
=
\frac{\delta_{30}}{1+\delta_{30}}
=
0.0002331873404817984\ldots,
\]

which is

\[
0.02331873404817984\ldots\%.
\]

\textbf{Proof.} Since \(R_*\ge R_K/(1+\delta_{30})\) and every rule in
the class has risk at least \(R_*\), the first inequality follows.
Subtract its ratio to \(R_K\) from one. \(\square\)

Therefore any competition-scale improvement necessarily violates at
least one assumption: nonnegativity, static/network independence, linear
cubature, limiting-kernel modeling, the node budget, or the original
function/kernel.

\subsection{3.3 MUB/Kerdock-line support
extremality}\label{mubkerdock-line-support-extremality}

Let \(B_1,\ldots,B_M\) be mutually unbiased orthonormal bases of
\(\mathbb R^d\), viewed as antipodal projective lines. Suppose an
antipodally symmetrized zonal kernel takes only three association values
on this universe:

\[
A=k(1),\qquad O=k(0),\qquad C=k(1/\sqrt d).
\]

Give line \((b,i)\) real weight \(w_{bi}\), with total mass one, and
define basis masses

\[
S_b=\sum_i w_{bi}.
\]

Up to a rule-independent constant, the kernel energy is

\[
R(w)= (O-C)\sum_b S_b^2 +(A-O)\sum_{b,i}w_{bi}^2.
\]

Define

\[
a=A-O,
\qquad
b=O-C.
\]

\textbf{Theorem 3.2 (association-scheme support extremality).} Assume

\[
a>0,\qquad b<0,\qquad a+bd>0.
\]

Among all real mass-one rules supported on at most \(P\) lines in the
fixed MUB universe, every minimum-energy rule:

\begin{enumerate}
\def\labelenumi{\arabic{enumi}.}
\tightlist
\item
  uses all \(P\) lines;
\item
  fills \(q=\lfloor P/d\rfloor\) complete bases;
\item
  uses at most one additional partial basis of size \(s=P-qd\);
\item
  assigns equal positive line weights within each active basis;
\item
  assigns positive basis masses proportional to \[
  \frac{1}{(O-C)+(A-O)/r_b}.
  \]
\end{enumerate}

\textbf{Proof.} If basis \(b\) contains \(r_b\) active lines,
Cauchy-Schwarz gives

\[
\sum_iw_{bi}^2\ge \frac{S_b^2}{r_b},
\]

with equality exactly for equal within-basis weights. Thus

\[
R(w)\ge \text{constant}+\sum_b c(r_b)S_b^2,
\qquad
c(r)=b+\frac ar.
\]

The assumptions imply \(c(r)>0\) for \(1\le r\le d\). Minimizing over
\((S_b)\) with \(\sum_bS_b=1\) gives

\[
S_b=\frac{c(r_b)^{-1}}{\sum_j c(r_j)^{-1}},
\qquad
R_{\min}=\text{constant}+\frac1{\sum_b h(r_b)},
\]

where

\[
h(r)=\frac1{c(r)}=\frac{r}{a+br},\qquad h(0)=0.
\]

On \([0,d]\),

\[
h'(r)=\frac{a}{(a+br)^2}>0,
\qquad
h''(r)=\frac{-2ab}{(a+br)^3}>0.
\]

Thus every available line is used and \(h\) is strictly convex. For
integers \(0<x\le y<d\), convexity implies

\[
h(x-1)+h(y+1)>h(x)+h(y).
\]

Repeatedly transferring one line from a smaller partial basis to a
larger one strictly increases \(\sum_bh(r_b)\), and therefore strictly
decreases risk, until only complete bases, at most one partial basis,
and empty bases remain. The formulas above force positive basis masses
and positive equal within-basis weights. \(\square\)

For the depth-32, \(d=256\) Kerdock universe, the certified signs are

\[
A-O>0,
\qquad
O-C<0,
\qquad
(A-O)+256(O-C)>0.
\]

Therefore the theorem applies to all static linear rules with arbitrary
real mass-one weights on the fixed 33,024 symmetrized lines. It does not
apply to nodes outside that universe, unpaired points, finite-width
adaptation, or nonlinear estimators.

\subsection{3.4 Why positive definiteness alone is
insufficient}\label{why-positive-definiteness-alone-is-insufficient}

The complete-basis conclusion does not follow from positive definiteness
alone. In \(d=4\), consider two real MUBs and

\[
k(t)=1+\lambda P_4(t),\qquad \lambda>0,
\]

where

\[
P_4(t)=\frac{16t^4-12t^2+1}{5}.
\]

This is positive definite, yet its association values satisfy \(O-C>0\).
Then \(h(r)=r/(a+br)\) is concave rather than convex, and a balanced
two-basis support beats one complete basis at a four-line budget. The
sign pattern, not positive definiteness alone, is the structural
mechanism.

\subsection{3.5 Arbitrary signed nodes: a stability lemma, not
closure}\label{arbitrary-signed-nodes-a-stability-lemma-not-closure}

Let an auxiliary minorant satisfy

\[
h(t)=\sum_{\ell=0}^Lc_\ell G_\ell(t),
\qquad c_\ell\ge0\ (\ell\ge1),
\qquad h(t)\le K(t),
\]

and put \(q=K-h\), \(q_1=q(1)\), \(M=\sup q\). For real weights summing
to one, define total negative mass

\[
\beta=\sum_{w_i<0}|w_i|.
\]

\textbf{Proposition 3.3.} For at most \(N\) nodes,

\[
E_K(w)
\ge
c_0+\frac{q_1}{N}-2M\beta(1+\beta).
\]

\textbf{Proof.} Positive definiteness of each Gegenbauer kernel gives

\[
\sum_{i,j}w_iw_jh(\langle x_i,x_j\rangle)\ge c_0.
\]

For the residual matrix, diagonal terms contribute \(q_1\sum_iw_i^2\);
same-sign off-diagonal terms are nonnegative; and the ordered absolute
mass of opposite-sign products is \(2\beta(1+\beta)\). Hence

\[
\sum_{i,j}w_iw_jq_{ij}
\ge
q_1\sum_iw_i^2-2M\beta(1+\beta).
\]

Finally, \(\sum_iw_i^2\ge1/N\). \(\square\)

The exact residual supremum in the audited certificate is

\[
M=\frac{156999263604490023}{9223372036854775808}
=0.017021894267861247\ldots.
\]

Numerically the lemma is too weak to close arbitrary signed-node
cubature: even a 10\% relative improvement requires only a tiny
certified lower bound on \(\beta\). This is a stability result, not a
universal signed-weight impossibility theorem.

\section{4. Information and observability
boundary}\label{information-and-observability-boundary}

\subsection{4.1 Runtime-information
value}\label{runtime-information-value}

\textbf{Theorem 4.1 (correction observability principle).} Let
\(e\in L^2(\Omega;H)\) and let \(\mathcal G\subseteq\mathcal F\) be all
information available to a runtime correction. Among all
\(\mathcal G\)-measurable corrections \(c\in L^2(\mathcal G;H)\),

\[
\inf_c\mathbb E\|e-c\|^2
=
\mathbb E\|e\|^2-
\mathbb E\|\mathbb E[e\mid\mathcal G]\|^2.
\]

The unique optimum is

\[
c^*=\mathbb E[e\mid\mathcal G].
\]

\textbf{Proof.} Conditional expectation is the orthogonal projection of
\(e\) onto the closed subspace \(L^2(\mathcal G;H)\). Put
\(m=\mathbb E[e\mid\mathcal G]\). For every admissible \(c\),

\[
e-c=(e-m)+(m-c),
\]

and the two terms are orthogonal. Therefore

\[
\mathbb E\|e-c\|^2
=
\mathbb E\|e-m\|^2+
\mathbb E\|m-c\|^2.
\]

The second term is minimized uniquely by \(c=m\), and the Pythagorean
identity gives the value. \(\square\)

Define the correction value

\[
V(\mathcal G;e)
=
\mathbb E\|\mathbb E[e\mid\mathcal G]\|^2.
\]

If \(\mathcal G_1\subseteq\mathcal G_2\), then

\[
V(\mathcal G_2;e)-V(\mathcal G_1;e)
=
\mathbb E\left\|
\mathbb E[e\mid\mathcal G_2]
-
\mathbb E[e\mid\mathcal G_1]
\right\|^2
\ge0.
\]

Thus replacing observations by a statistic cannot increase their
correction value.

\subsection{4.2 Finite-dictionary adaptive coefficient
theorem}\label{finite-dictionary-adaptive-coefficient-theorem}

Let \(A(\omega):\mathbb R^m\to H\) be a random linear feature operator.
A coefficient vector \(a\), measurable with respect to runtime
information \(\mathcal G\), produces correction \(Aa\).

Define conditional normal-equation quantities

\[
\Sigma_{\mathcal G}
=
\mathbb E[A^*A\mid\mathcal G],
\qquad
\mu_{\mathcal G}
=
\mathbb E[A^*e\mid\mathcal G].
\]

\textbf{Theorem 4.2 (adaptive dictionary value).} The optimal
coefficient is

\[
a^*_{\mathcal G}
=
\Sigma_{\mathcal G}^{\dagger}\mu_{\mathcal G},
\]

with minimum-norm convention, and the attainable gain is

\[
V_A(\mathcal G;e)
=
\mathbb E\left[
\mu_{\mathcal G}^{\mathsf T}
\Sigma_{\mathcal G}^{\dagger}
\mu_{\mathcal G}
\right].
\]

\textbf{Proof.} Condition on \(\mathcal G\). The excess risk relative to
zero correction is

\[
a^{\mathsf T}\Sigma_{\mathcal G}a
-2a^{\mathsf T}\mu_{\mathcal G}.
\]

If \(v\in\ker\Sigma_{\mathcal G}\), then

\[
0=v^{\mathsf T}\Sigma_{\mathcal G}v
=
\mathbb E[\|Av\|^2\mid\mathcal G],
\]

so \(Av=0\) almost surely conditionally and therefore
\(v^{\mathsf T}\mu_{\mathcal G}=0\). Hence
\(\mu_{\mathcal G}\in\operatorname{range}\Sigma_{\mathcal G}\), and
completing the square with the Moore-Penrose inverse gives the stated
optimizer and gain. \(\square\)

If the full per-network oracle observes \((A,e)\), then

\[
a_{\rm oracle}= (A^*A)^\dagger A^*e,
\]

and its gain is

\[
V_A({\rm oracle};e)
=
\mathbb E\|P_{\operatorname{range}(A)}e\|^2.
\]

This separates two questions:

\begin{itemize}
\tightlist
\item
  \textbf{Capacity:} does the dictionary span useful corrections?
\item
  \textbf{Observability:} can legal runtime information recover their
  signed coefficients?
\end{itemize}

For nested runtime and oracle classes, define

\[
\operatorname{ObservabilityRatio}
=
\frac{V_A(\mathcal G_{
m runtime};e)}
{V_A(\mathcal G_{
m oracle};e)}
\in[0,1].
\]

A large oracle gain with a small observability ratio is exactly the
``capacity without transferable phase'' pattern seen in the signed-probe
experiments.

\subsection{4.3 Phase-flip impossibility
theorem}\label{phase-flip-impossibility-theorem}

\textbf{Theorem 4.3.} Suppose \(\tau:\Omega\to\Omega\) is a
measure-preserving involution such that

\begin{enumerate}
\def\labelenumi{\arabic{enumi}.}
\tightlist
\item
  every runtime observable in \(\mathcal G\) is invariant under
  \(\tau\);
\item
  the dictionary cross-moment flips sign: \[
  A(\tau\omega)^*e(\tau\omega)
  =-A(\omega)^*e(\omega).
  \]
\end{enumerate}

Then

\[
\mu_{\mathcal G}=0
\]

and every \(\mathcal G\)-measurable dictionary coefficient rule has
nonnegative excess risk. The unique minimum-norm optimum is zero
correction.

\textbf{Proof.} Invariance of \(\mathcal G\), measure preservation, and
anti-invariance of \(A^*e\) imply

\[
\mathbb E[A^*e\mid\mathcal G]
=
-\mathbb E[A^*e\mid\mathcal G],
\]

so it is zero. Theorem 4.2 then gives zero attainable gain. \(\square\)

The same conclusion holds if the conditional law of \(A^*e\) given
\(\mathcal G\) is centrally symmetric. This is the clean theorem one
would need to prove for a specific weight-summary representation in
order to convert empirical phase reversal into a population
impossibility result.

\subsection{4.4 Haar relative-orientation
theorem}\label{haar-relative-orientation-theorem}

Let \(G=O(d)\) act transitively on \(S^{d-1}\), with normalized Haar
measure \(dg\). Let \(Q=\sum_iw_i\delta_{x_i}\) have total mass one, and
define its rotation by

\[
Q_gf=\sum_iw_i f(gx_i).
\]

\textbf{Theorem 4.4 (orientation-blind no-value theorem).} For every
integrable Hilbert-valued \(f\),

\[
\int_G Q_gf\,dg=I(f).
\]

Consequently, if \(g\) is Haar and independent of all orientation-blind
runtime information \(\mathcal G\subseteq\sigma(f)\), then

\[
\mathbb E[Q_gf-I(f)\mid\mathcal G]=0.
\]

No \(\mathcal G\)-measurable correction can lower mean squared error.

\textbf{Proof.} For every fixed node \(x_i\), the pushforward of Haar
measure under \(g\mapsto gx_i\) is uniform spherical measure. Therefore

\[
\int_G f(gx_i)\,dg=I(f).
\]

Summing with \(\sum_iw_i=1\) proves the first identity. Conditional mean
zero follows because \(g\) is independent of \(\mathcal G\), and Theorem
4.1 gives zero correction value. \(\square\)

This theorem rigorously closes corrections based only on
rotation-invariant network summaries when relative design orientation is
independently Haar randomized. It does \textbf{not} close
orientation-sensitive white-box methods, network-derived nodes, or
corrections using design evaluations.

\subsection{4.5 Same-design common-bias
non-identifiability}\label{same-design-common-bias-non-identifiability}

Suppose same-design estimates obey

\[
Z_i=\mu+b+\varepsilon_i,
\qquad i=1,\ldots,k,
\]

where the joint noise law does not depend on \((\mu,b)\).

\textbf{Theorem 4.5.} The transformation

\[
(\mu,b)\mapsto(\mu+t,b-t)
\]

leaves the complete observation law unchanged. For every estimator
\(T(Z)\) of \(\mu\), at least one of the two gauge-equivalent parameter
points has squared-error risk at least

\[
\frac{\|t\|^2}{4}.
\]

\textbf{Proof.} The observation law depends on \((\mu,b)\) only through
\(\mu+b\). Pointwise, for any estimate \(x\),

\[
\|x-\mu\|^2+
\|x-(\mu+t)\|^2
=
2\left\|x-\mu-\frac t2\right\|^2+
\frac{\|t\|^2}{2}
\ge
\frac{\|t\|^2}{2}.
\]

Taking expectations under the common observation law shows that at least
one risk is at least \(\|t\|^2/4\). \(\square\)

Centered diagnostics

\[
D_i=Z_i-\overline Z
\]

depend only on centered noise. They can estimate dispersion or
instability, but not the shared absolute defect in this model.

The theorem is model-specific. It fails with unequal known loadings, for
example

\[
Z_1=\mu+b+\varepsilon_1,
\qquad
Z_2=\mu-b+\varepsilon_2,
\]

because \(Z_1-Z_2\) directly contains \(2b\).

\subsection{4.6 General indistinguishability lower
bound}\label{general-indistinguishability-lower-bound}

\textbf{Theorem 4.6.} Suppose two admissible worlds
\(\theta_0,\theta_1\) induce exactly the same runtime transcript
distribution but have targets \(y_0,y_1\in H\). Then every estimator
\(T\) based on that transcript satisfies

\[
\max_{j\in\{0,1\}}
\mathbb E_j\|T-y_j\|^2
\ge
\frac{\|y_1-y_0\|^2}{4}.
\]

\textbf{Proof.} The transcript and hence \(T\) have the same
distribution under both worlds. Apply the parallelogram calculation from
Theorem 4.5 with \(t=y_1-y_0\). \(\square\)

This theorem converts a concrete indistinguishable-pair construction
into a query or statistic lower bound. No such pair has yet been
constructed for the full legal width-256 white-box network class.

\subsection{4.7 Why a universal information impossibility theorem is
false}\label{why-a-universal-information-impossibility-theorem-is-false}

\textbf{Proposition 4.7 (full-information obstruction).} If the complete
network weights \(W\) determine the target and baseline error, then

\[
\mathbb E[e\mid W]=e,
\qquad
V(\sigma(W);e)=\mathbb E\|e\|^2.
\]

Thus full weights contain, information-theoretically, enough information
for a perfect correction.

\textbf{Proof.} The error is \(\sigma(W)\)-measurable, so its
conditional expectation given \(W\) is itself. Substitute in Theorem
4.1. \(\square\)

Therefore any meaningful lower bound must restrict at least one of:

\begin{itemize}
\tightlist
\item
  the statistics extracted from the weights;
\item
  the number or geometry of network evaluations;
\item
  the estimator family;
\item
  the permitted arithmetic or time complexity;
\item
  the information protocol.
\end{itemize}

The experiments alone cannot prove that no full-white-box algorithm
exists.

\section{5. Correction and downstream-replay
mathematics}\label{correction-and-downstream-replay-mathematics}

\subsection{5.1 Exact risk identity}\label{exact-risk-identity}

For a proposed direction \(u\in L^2(\Omega;H)\) and scalar \(\alpha\),
define

\[
R(\alpha)=\mathbb E\|e-\alpha u\|^2.
\]

Let

\[
R_0=\mathbb E\|e\|^2,
\quad
C=\mathbb E\langle e,u\rangle,
\quad
U=\mathbb E\|u\|^2.
\]

\textbf{Theorem 5.1.}

\[
R(\alpha)=R_0-2\alpha C+\alpha^2U.
\]

For \(U>0\),

\[
\alpha_*=\frac CU,
\qquad
R(\alpha_*)=R_0-\frac{C^2}{U}.
\]

If only \(\alpha\ge0\) is legal, replace \(C\) by \((C)_+\).

\textbf{Proof.} Expand the squared norm and minimize the resulting
quadratic. \(\square\)

Magnitude, disagreement, and oracle span are insufficient. The binding
quantity is the signed error-correction inner product.

\subsection{5.2 Replacement in a correctable
subspace}\label{replacement-in-a-correctable-subspace}

Let \(\mathcal S\subset L^2(\Omega;H)\) be a closed linear subspace.
Decompose

\[
s=P_{\mathcal S}e,
\qquad
r=e-s,
\qquad
r\perp\mathcal S.
\]

Suppose an anchor estimates \(s\) as \(\widehat s=s+n\).

\textbf{Theorem 5.2.} If \(n\in\mathcal S\), then

\[
\mathbb E\|e-\widehat s\|^2
=
\mathbb E\|r\|^2+
\mathbb E\|n\|^2.
\]

Full replacement improves exactly when

\[
\mathbb E\|n\|^2<\mathbb E\|s\|^2.
\]

\textbf{Proof.} Since \(e-\widehat s=r-n\) and \(r\perp n\), Pythagoras
gives the identity. Compare with
\(\mathbb E\|e\|^2=\mathbb E\|r\|^2+\mathbb E\|s\|^2\). \(\square\)

If \(n\) is not in \(\mathcal S\), the exact formula is

\[
\mathbb E\|e-\widehat s\|^2
=
\mathbb E\|r\|^2+
\mathbb E\|n\|^2
-2\mathbb E\langle r,n\rangle.
\]

Thus a norm-only replacement gate requires the subspace or orthogonality
assumption.

\subsection{5.3 Downstream-weighted anchor
criterion}\label{downstream-weighted-anchor-criterion}

Condition on a realized network and protected particle cloud. Let

\[
d=\mu_{31}^{K}-\mu_{31}^*
\]

be the protected layer-31 mean defect and

\[
\xi=\widehat\mu_{31}-\mu_{31}^*
\]

be anchor error. Let \(J\) be the frozen linearized map from layer-31
mean perturbations to scored final outputs. The relevant errors are

\[
s=Jd,
\qquad
n=J\xi.
\]

Under Theorem 5.2's subspace assumptions, full replacement improves
exactly when

\[
\eta_J^2
:=
\frac{\mathbb E\|J\xi\|^2}
{\mathbb E\|Jd\|^2}
<1.
\]

No universal threshold in unweighted layer-31 Euclidean error exists.
For a unit direction \(v\), the linearized break-even amplitude is

\[
\epsilon_*(v)
=
\frac{\|Jd\|}{\|Jv\|},
\]

which varies with direction unless \(J^*J\) is scalar on the tested
span.

This formally justifies rescoring anchors by exact downstream replay
rather than by a scalar coefficient-space threshold.

\subsection{5.4 ReLU crossing remainder}\label{relu-crossing-remainder}

For \(\phi(z)=\max(z,0)\), define

\[
r(z,t)=\phi(z+t)-\phi(z)-\mathbf 1_{\{z>0\}}t.
\]

\textbf{Lemma 5.3.} For all \(z,t\in\mathbb R\),

\[
|r(z,t)|
\le
|t|\mathbf 1_{\{|z|\le|t|\}}.
\]

\textbf{Proof.} If \(z\) and \(z+t\) have the same sign, ReLU is affine
on the segment and the remainder is zero. A sign change implies
\(|z|\le|t|\), and direct inspection bounds the discrepancy by \(|t|\).
\(\square\)

If the conditional density of \(z\) is bounded by \(L\) near zero, then

\[
\mathbb E[r(z,t)^2\mid t]
\le
|t|^2\Pr(|z|\le|t|\mid t)
\le
2L|t|^3.
\]

For a vector replay with linearized error \(r_0-n\) and nonlinear
remainder \(q\), Minkowski gives

\[
\mathbb E\|r_0-n+q\|^2
\le
\left(
\sqrt{\mathbb E\|r_0-n\|^2}
+
\sqrt{\mathbb E\|q\|^2}
\right)^2.
\]

A sufficient exact-nonlinear improvement condition is that the
right-hand side be below baseline risk.

\section{6. Exact control mathematics}\label{exact-control-mathematics}

\subsection{6.1 Uniform control-annihilation
theorem}\label{uniform-control-annihilation-theorem}

For a fixed quadrature rule \(Q\), define its exactness space

\[
\mathcal N_Q=\{g:I(g)=Q(g)\}.
\]

\textbf{Theorem 6.1.} For every \(f\), every \(g\in\mathcal N_Q\), and
every scalar \(\alpha\),

\[
I(\alpha g)+Q(f-\alpha g)=Q(f).
\]

More generally, suppose \(\{g_\eta:\eta\in\Xi\}\subseteq\mathcal N_Q\)
uniformly. The parameter \(\widehat\eta\) may be selected adaptively
from the network, pilots, or quadrature outputs. If the same selected
function is used in both \(I(g_{\widehat\eta})\) and
\(Q(g_{\widehat\eta})\), the estimator remains exactly \(Q(f)\)
pathwise.

\textbf{Proof.} Linearity gives

\[
I(\alpha g)+Q(f-\alpha g)
=Q(f)+\alpha(I(g)-Q(g))
=Q(f).
\]

The adaptive statement is pointwise in the selected \(\widehat\eta\).
\(\square\)

Thus increasing model expressivity within a uniformly annihilated family
cannot help. This closes named low-degree and exact ReLU-Stein families,
not all high-degree or node-dependent controls.

\subsection{6.2 Symmetrized spherical Poisson
controls}\label{symmetrized-spherical-poisson-controls}

For \(u,x\in S^{d-1}\) and \(|r|<1\), define

\[
P_r(\langle u,x\rangle)
=
\frac{1-r^2}
{(1-2r\langle u,x\rangle+r^2)^{d/2}}.
\]

\textbf{Theorem 6.2.} With normalized spherical measure \(\sigma\),

\[
\int_{S^{d-1}}P_r(\langle u,x\rangle)\,d\sigma(x)=1.
\]

Therefore the antipodally symmetrized control

\[
P_r^{\rm even}(t)
=
\frac12\bigl(P_r(t)+P_r(-t)\bigr)
\]

also has exact mean one.

\textbf{Proof.} \(P_r\) is the Poisson kernel for the unit ball. Its
integral against boundary measure is the harmonic extension of the
constant boundary function one, which is identically one. Equivalently,
expand in zonal spherical harmonics: the degree-zero coefficient is one
and every positive-degree harmonic integrates to zero. \(\square\)

The zonal expansion has the form

\[
P_r(t)=\sum_{\ell=0}^{\infty}r^\ell Z_\ell(t),
\]

where \(Z_\ell\) is the degree-\(\ell\) zonal reproducing kernel.
Symmetrization cancels odd degrees:

\[
P_r^{\rm even}(t)
=
\sum_{k=0}^{\infty}r^{2k}Z_{2k}(t).
\]

For \(r\ne0\), every even degree has a nonzero coefficient. Hence
analytically integrable does not imply low harmonic degree.

\subsection{6.3 Closed-form spherical mean of a biased projected
ReLU}\label{closed-form-spherical-mean-of-a-biased-projected-relu}

Let \(X\) be uniform on the radius-\(R\) sphere in \(\mathbb R^d\). For
\(a\in\mathbb R^d\), put

\[
\rho=R\|a\|,
\qquad
T=\frac{\langle a,X\rangle}{\rho}
\]

when \(\rho>0\). Then \(T\in[-1,1]\) has density

\[
f_d(t)=c_d(1-t^2)^{(d-3)/2},
\qquad
c_d=
\frac{\Gamma(d/2)}
{\sqrt\pi\,\Gamma((d-1)/2)}.
\]

Consider

\[
g_{a,b}(X)=(\langle a,X\rangle+b)_+.
\]

\textbf{Theorem 6.3.} If \(\rho=0\), then \(\mathbb E g_{a,b}=b_+\). If
\(b\ge\rho\), then \(\mathbb E g_{a,b}=b\); if \(b\le-\rho\), it is
zero. For \(|b|<\rho\), put

\[
s=-\frac b\rho,
\qquad
\beta=\frac{d-1}{2}.
\]

Then

\[
\boxed{
\mathbb E g_{a,b}
=
\frac{\rho c_d}{d-1}(1-s^2)^\beta
+
 b\left[
\frac12-
\frac{\operatorname{sgn}(s)}2
I_{s^2}\left(\frac12,\beta\right)
\right]
}
\]

where \(I_x(p,q)\) is the regularized incomplete beta function, with
\(\operatorname{sgn}(0)=0\).

\textbf{Proof.} The active region is \(T>s\). Therefore

\[
\mathbb E g_{a,b}
=
\rho c_d\int_s^1t(1-t^2)^{(d-3)/2}\,dt
+b\Pr(T>s).
\]

The first integral is

\[
\int_s^1t(1-t^2)^{(d-3)/2}\,dt
=
\frac{(1-s^2)^{(d-1)/2}}{d-1}.
\]

Symmetry of \(T\) and the substitution \(u=t^2\) give

\[
\Pr(T>s)
=
\frac12-
\frac{\operatorname{sgn}(s)}2
I_{s^2}\left(\frac12,\frac{d-1}{2}\right).
\]

Substitution yields the formula. \(\square\)

For the experiment's normalized direction rows \(\|a\|=1\), use
\(\rho=R\). The verifier checks the formula against adaptive numerical
integration to approximately \texttt{1e-14} absolute error. The
experiment used 384-point Gauss-Jacobi quadrature; the closed form
removes that numerical dependency for future implementations.

\section{7. Replication and compute
economics}\label{replication-and-compute-economics}

Suppose \(m\) estimators have errors \(e_1,\ldots,e_m\in H\), each with
risk \(R_0\), and pairwise normalized covariance

\[
\mathbb E\langle e_i,e_j\rangle=\rho R_0
\qquad(i\ne j).
\]

Their average has risk

\[
\mathbb E\left\|\frac1m\sum_{i=1}^me_i\right\|^2
=
R_0\left(\rho+\frac{1-\rho}{m}\right).
\]

\textbf{Proof.} Expand the squared norm. There are \(m\) diagonal terms
and \(m(m-1)\) ordered off-diagonal terms. Divide by \(m^2\).
\(\square\)

If compute grows by a factor \(c_m\), the adjusted ratio is

\[
c_m\left(\rho+\frac{1-\rho}{m}\right).
\]

For independent full-cost replicas, \(\rho=0\) and \(c_m=m\), so the
adjusted ratio is exactly one. Ordinary independent replication cannot
improve an MSE-times-linear-compute score. It can win only through
sublinear shared arithmetic, negative covariance, unequal quality/cost
allocation, or a different scoring rule.

\section{8. Finite-width and transformed-residual
boundaries}\label{finite-width-and-transformed-residual-boundaries}

\subsection{8.1 Kernel perturbation
theorem}\label{kernel-perturbation-theorem}

Let \(P\) be the target distribution and \(Q=\sum_iw_i\delta_{x_i}\).
Put \(\nu_Q=P-Q\). For kernels \(K\) and \(\widetilde K\), define

\[
R_K(Q)=\iint K(x,y)\,d\nu_Q(x)d\nu_Q(y).
\]

\textbf{Theorem 8.1.} If

\[
\|K-\widetilde K\|_\infty\le\varepsilon,
\qquad
\sum_i|w_i|\le B,
\]

then

\[
|R_K(Q)-R_{\widetilde K}(Q)|
\le
\varepsilon(1+B)^2.
\]

\textbf{Proof.} \(\|\nu_Q\|_{\rm TV}\le1+B\), so

\[
\left|\iint(K-\widetilde K)\,d\nu_Qd\nu_Q\right|
\le
\|K-\widetilde K\|_\infty\|\nu_Q\|_{\rm TV}^2.
\]

\(\square\)

If \(Q_K\) has additive suboptimality at most \(g\) for \(K\), then

\[
R_{\widetilde K}(Q_K)-\inf_QR_{\widetilde K}(Q)
\le
 g+2\varepsilon(1+B)^2.
\]

For nonnegative mass-one rules, \(B=1\), so the perturbation allowance
is \(g+8\varepsilon\).

The certified limiting-kernel gap is extremely small. A qualitative
statement that width-256 kernels ``approach'' the limiting kernel is
insufficient; a quantitative error bound at the relevant scale is
required.

\subsection{8.2 Optimizer instability under tiny PSD
perturbations}\label{optimizer-instability-under-tiny-psd-perturbations}

Let two rules have linearly independent error functionals \(L_0,L_1\).
Choose bounded \(\phi,\psi\) satisfying

\[
L_0\phi=0,
\quad L_1\phi=1,
\qquad
L_0\psi=1,
\quad L_1\psi=0.
\]

Set

\[
K_\infty=g\,\phi\otimes\phi,
\qquad
K_m=K_\infty+\varepsilon\,\psi\otimes\psi.
\]

Both kernels are positive semidefinite. Under \(K_\infty\), rule zero
has risk zero and rule one has risk \(g\). Under \(K_m\), their risks
are \(\varepsilon\) and \(g\). For \(\varepsilon>g\), the optimizer
reverses. Since \(g\) can be arbitrarily small, optimizer identity is
not structurally stable without a quantitative gap-versus-perturbation
comparison.

\subsection{8.3 Residual-kernel
recertification}\label{residual-kernel-recertification}

Let \(g_\theta\) be a network-dependent analytic control with exactly
available integral. Define

\[
h_\theta=f_\theta-g_\theta
\]

and estimator

\[
\widehat I_Q(f_\theta)=I(g_\theta)+Q(h_\theta).
\]

Then

\[
\widehat I_Q(f_\theta)-I(f_\theta)
=(Q-I)h_\theta.
\]

If

\[
K_{\rm res}(x,y)
=
\mathbb E[h_\theta(x)h_\theta(y)],
\]

then

\[
\mathbb E|
\widehat I_Q(f_\theta)-I(f_\theta)
|^2
=
\iint K_{\rm res}(x,y)\,d\nu_Q(x)d\nu_Q(y).
\]

Thus analytic-plus-residual estimation returns to a kernel-energy
problem, but for \(K_{\rm res}\), not the original deep-ReLU kernel.
Kerdock must be recertified after the transformation. Original
optimality does not automatically persist.

\section{9. What the frozen experiments
prove}\label{what-the-frozen-experiments-prove}

The deterministic verifier bundled with this document checks all 154
entries in the reopened-path SHA-256 manifest and recomputes the
principal terminal metrics from the saved row-level arrays.

The finite statements are:

\begin{enumerate}
\def\labelenumi{\arabic{enumi}.}
\tightlist
\item
  The frozen Poisson terminal candidate has raw ratio \[
  1.0379391539423775>1.
  \]
\item
  The frozen projected-ReLU candidate over 48 terminal networks has raw
  ratio \[
  1.012682140618133>1,
  \] and favorable adjusted ratio \[
  1.0668717656262163>1.
  \]
\item
  The frozen network-derived signed-probe rule has raw ratio \[
  1.5573136434384325>1.
  \]
\item
  The same signed dictionary has a per-network ridge oracle ratio \[
  0.13117636626944396,
  \] below the \texttt{4.34x} target ratio, proving that its span has
  enough same-network capacity while its tested transferable coefficient
  rule fails.
\item
  The exact Poisson mean-one identity and the closed-form projected-ReLU
  mean agree with independent numerical integration to near machine
  precision.
\end{enumerate}

These are propositions about the frozen files, candidate definitions,
and IEEE-754 computations. They do not establish that the synthetic
network generator equals the protected competition distribution or that
every member of a broader estimator class fails.

\subsection{9.1 Empirical observability
diagnostics}\label{empirical-observability-diagnostics}

A post-hoc decomposition, explicitly not a generalization claim, found:

\begin{longtable}[]{@{}
  >{\raggedright\arraybackslash}p{(\columnwidth - 6\tabcolsep) * \real{0.2000}}
  >{\raggedleft\arraybackslash}p{(\columnwidth - 6\tabcolsep) * \real{0.2667}}
  >{\raggedleft\arraybackslash}p{(\columnwidth - 6\tabcolsep) * \real{0.2667}}
  >{\raggedleft\arraybackslash}p{(\columnwidth - 6\tabcolsep) * \real{0.2667}}@{}}
\toprule\noalign{}
\begin{minipage}[b]{\linewidth}\raggedright
Family
\end{minipage} & \begin{minipage}[b]{\linewidth}\raggedleft
Post-hoc global ratio
\end{minipage} & \begin{minipage}[b]{\linewidth}\raggedleft
Per-network oracle ratio
\end{minipage} & \begin{minipage}[b]{\linewidth}\raggedleft
Empirical global/oracle gain ratio
\end{minipage} \\
\midrule\noalign{}
\endhead
\bottomrule\noalign{}
\endlastfoot
Network-derived signed probes & \texttt{0.93821} & \texttt{0.13118} &
\texttt{0.07112} \\
Random signed probes & \texttt{0.95802} & \texttt{0.33681} &
\texttt{0.06330} \\
Fixed projected-ReLU direction with scalar coefficient &
\texttt{0.99897} & \texttt{0.84762} & \texttt{0.00676} \\
Fixed Poisson direction with scalar coefficient & \texttt{0.99210} &
\texttt{0.78358} & \texttt{0.03652} \\
\end{longtable}

For the network-derived signed probes, leave-one-network-out fitting was
harmful (\texttt{1.09909x}). These diagnostics strongly motivate a
phase-observability theorem, but they do not themselves prove one. The
required population proof obligation is Theorem 4.3: identify an exact
conditional symmetry or bound the conditional cross-moment.

\section{10. Claims that remain open}\label{claims-that-remain-open}

The following are not proved by this package:

\begin{enumerate}
\def\labelenumi{\arabic{enumi}.}
\tightlist
\item
  Exact width-256 optimality of Kerdock.
\item
  A universal lower bound against all algorithms using the complete
  network weights.
\item
  Optimality over arbitrary signed spherical nodes.
\item
  Impossibility of network-adaptive support or weights.
\item
  Impossibility of nonlinear postprocessing.
\item
  Impossibility of all biased, deep, nonhomogeneous, or nonpolynomial
  analytic controls.
\item
  Population-level zero observability for the exact signed-probe
  representation.
\item
  Equality between the synthetic seed distribution and the protected
  competition cohort.
\item
  A computational lower bound for evaluating the target from full
  weights.
\end{enumerate}

Any manuscript sentence equivalent to ``no statistical path exists'' is
therefore false or unsupported.

\section{11. Strongest defensible integrated
theorem}\label{strongest-defensible-integrated-theorem}

The results can be stated compactly as follows.

\textbf{Integrated two-boundary theorem.} For the dimension-256,
depth-32 limiting ReLU kernel:

\begin{enumerate}
\def\labelenumi{\arabic{enumi}.}
\tightlist
\item
  complete Kerdock is computer-assisted certified within
  \texttt{0.023324172950039\%} of the static nonnegative optimum at the
  stated node budget;
\item
  inside the fixed 33,024-line Kerdock universe, complete bases plus at
  most one partial basis exactly optimize all static linear mass-one
  rules, even with arbitrary real weights;
\item
  for any explicitly specified runtime information class, attainable
  adaptive correction gain is exactly the squared Hilbert projection of
  the baseline error onto that information class;
\item
  equal-loading common-bias information, independent
  Haar-orientation-blind information, and feature classes satisfying a
  phase-flip symmetry have zero absolute-phase correction value;
\item
  full network weights do not satisfy a generic
  information-impossibility premise, so broader white-box and
  finite-width methods remain mathematically open.
\end{enumerate}

Therefore a material winner must introduce either:

\begin{itemize}
\tightlist
\item
  support/weights outside the certified geometric class;
\item
  information that breaks the absolute-phase symmetry;
\item
  a transformed residual whose kernel is materially easier;
\item
  finite-width-specific analytic structure;
\item
  or a computational mechanism not represented by the tested correction
  families.
\end{itemize}

\section{12. Reproducibility and trust
base}\label{reproducibility-and-trust-base}

Bundled files:

\begin{itemize}
\tightlist
\item
  \texttt{WHESTBENCH\_COMPLETE\_PROOF\_PACKAGE\_20260730.md} - this
  proof document;
\item
  \texttt{WHESTBENCH\_COMPLETE\_PROOF\_PACKAGE\_20260730.tex} - LaTeX
  source;
\item
  \texttt{WHESTBENCH\_COMPLETE\_PROOF\_PACKAGE\_20260730.pdf} - rendered
  proof document;
\item
  \texttt{verify\_complete\_proof\_package.py} - deterministic
  frozen-artifact verifier;
\item
  \texttt{WHESTBENCH\_PROOF\_CERTIFICATE\_20260730.json} - verifier
  output;
\item
  \texttt{OBSERVABILITY\_EMPIRICAL\_AUDIT\_20260730.json} - explicitly
  post-hoc empirical diagnostic;
\item
  \texttt{reopened\_paths\_repro\_20260730/} - scripts, frozen arrays,
  split registry, preregistrations, hashes, and results;
\item
  \texttt{sources/} - materialized theorem-source summaries and the T16
  primal-dual certificate.
\end{itemize}

The proof dependencies T22/T30 remain computer-assisted, not
proof-assistant formalized. Before public release, a qualified human
should review every analytic argument, rerun the original interval
engines and CI matrix, publish an externally authenticated archive
digest, and disclose the extent of language-model assistance.

\end{document}
